\chapter{\texttt{core.numerical} --- general numerical propagation}
\label{ch:core-numerical}

This chapter covers the medium-independent numerical evolutor: exact
matrix-exponential integration of the propagation Hamiltonian over a
piecewise-constant-density discretisation of the path. It makes no
assumption on the size of matter effects and works unmodified for the
Standard Model, NSI, the 3+1 sterile extension, or any combination
(\cref{ch:core-common,ch:core-sm,ch:core-bsm}), since it consumes
\pyfunc{hamiltonian\_flavour} directly rather than any scenario-specific
formula. It is the general-purpose reference method against which the
perturbative approximation (\cref{ch:core-perturbative}) is validated.

\section{\texttt{geometry.py} --- trajectory discretisation}

\modulehead{core/numerical/geometry.py}{\pyclass{Trajectory}: the geometry
container passed from a medium profile to the numerical evolutor, and
\pyfunc{segment\_sample\_points}, the quadrature-rule helper that builds
it.}

\begin{theorybox}
Composing segment evolutors (\cref{ch:core-common}, \texttt{evolutor.py}
section) requires splitting a propagation path into $N$ segments delimited
by $N+1$ boundary points along some geometry coordinate $x$ (a chord
length, a radius, a zenith-angle parametrisation -- whatever the medium
finds natural). For each segment, the numerical core needs exactly two
further pieces of information, both computed once by the medium-specific
module that builds the \pyclass{Trajectory} and never revisited by
\texttt{core.numerical} itself:
\begin{itemize}
  \item the \textbf{dimensionless evolution length} of the segment,
    $\Delta x^{\rm evo}_j = \Delta x_j / L_{\rm scale}$, i.e.\ the physical
    segment length divided by the same \code{evolution\_scale\_m} used to
    non-dimensionalise the Hamiltonian (\cref{ch:core-common},
    \texttt{potential.py} section) -- this is what actually multiplies
    $H$ in the matrix exponential below, not $\Delta x_j$ itself;
  \item a \textbf{sample coordinate} $\bar x_j$ per segment, at which the
    medium's electron density (or any other property entering $H$) is
    evaluated once and treated as \emph{constant} over the whole segment
    -- the piecewise-constant-density approximation that makes the
    per-segment problem exactly solvable (next section).
\end{itemize}
\pyfunc{segment\_sample\_points} implements three quadrature rules for
$\bar x_j$ from the $N+1$ boundary coordinates $x_0,\dots,x_N$:
\begin{equation}
  \bar x_j =
  \begin{cases}
    \tfrac12\bigl(x_j+x_{j+1}\bigr) & \text{\code{"midpoint"}\ (default)}, \\
    x_j & \text{\code{"left"}}, \\
    x_{j+1} & \text{\code{"right"}}.
  \end{cases}
\end{equation}
\keyresult{The midpoint rule is second-order accurate} in the segment
width for a smoothly varying density (the segment's leading-order
mismatch between the true, continuously varying $n_e(x)$ and the
constant $n_e(\bar x_j)$ used on it cancels to first order when $\bar x_j$
is the segment centre, by the same Taylor-expansion argument as the
midpoint rule in quadrature); the endpoint rules (\code{"left"}/
\code{"right"}) are only first-order accurate, and exist mainly to match
legacy or directionally-biased sampling conventions.
\end{theorybox}

\begin{implbox}
\begin{itemize}
  \item \pyclass{Trajectory} (dataclass): \code{x} (shape $(N{+}1,)$, the
    geometry-coordinate boundaries), \code{dx\_evolution} (shape $(N,)$,
    the dimensionless evolution lengths $\Delta x^{\rm evo}_j$ above),
    \code{sample\_x} (shape $(N,)$, the $\bar x_j$ above), and \code{meta}
    (a plain \code{dict} of medium-specific bookkeeping -- layer indices,
    crossing flags, ... -- that \texttt{core.numerical} never interprets).
  \item \pyfunc{segment\_sample\_points(x, method)}: builds \code{sample\_x}
    from \code{x} via the three rules above; broadcasts over any leading
    batch dimensions, and returns an empty (zero-segment) result for a
    single-boundary-point input.
  \item Every medium-specific geometry module (e.g.\
    \texttt{medium.earth.geometry}, \texttt{medium.atmosphere.profile})
    builds its own \code{x} in whatever coordinate is natural for that
    medium, converts it to \code{dx\_evolution} using \emph{that medium's}
    evolution scale, and calls \pyfunc{segment\_sample\_points} to fill
    \code{sample\_x} -- \pyclass{Trajectory} is the single, medium-agnostic
    hand-off point to \texttt{core.numerical}.
\end{itemize}
\end{implbox}

\section{\texttt{evolutor.py} --- matrix-exponential composition}

\modulehead{core/numerical/evolutor.py}{\pyfunc{evolutor\_numerical\_segment}
and \pyfunc{evolutor\_numerical}: exact per-segment matrix exponentials,
composed into the total evolutor.}

\begin{theorybox}
\subsection*{Exact solution on a piecewise-constant-density segment}
The propagation Hamiltonian $H(x)$ varies continuously with $x$ through the
local electron density $n_e(x)$. Given the \pyclass{Trajectory} of the
previous section, \texttt{core.numerical} approximates $H(x)$ as
\textbf{piecewise constant}: on segment $j$ (between boundaries $x_j$ and
$x_{j+1}$), $H(x)\approx H_j$, where
\begin{equation}
  H_j \equiv \text{\pyfunc{hamiltonian\_flavour}}\bigl(\text{oscillation},\,E,\,n_e(\bar x_j)\bigr)
\end{equation}
is the full flavour-basis Hamiltonian (\cref{ch:core-common}) evaluated
once at the segment's sample point $\bar x_j$. \textbf{This is the
\emph{only} approximation in the method}: given a constant $H_j$, the
position-dependent Schr\"odinger equation
$i\,\mathrm{d}S/\mathrm{d}x=H_j\,S$ (\cref{ch:core-common},
\texttt{evolutor.py} section) integrates \emph{exactly} -- not
approximately -- over the segment, via the matrix exponential:
\begin{equation}
  \keyresult{S_j = \exp\bigl(-i\,H_j\,\Delta x^{\rm evo}_j\bigr)}.
\end{equation}
No eigendecomposition, diagonalisation, or perturbative expansion is
involved: \pyfunc{torch.linalg.matrix\_exp} evaluates this directly (via a
scaling-and-squaring / Pad\'e algorithm), batched over every segment and
every other broadcast dimension (energy, angle, ...) simultaneously, and
remains differentiable end-to-end. Because $H_j$ is built with
\pyfunc{hamiltonian\_flavour} -- the single, fully general
Hamiltonian-assembly function (\cref{ch:core-common,ch:core-sm,ch:core-bsm}) --
this formula is correct as written for the 3-flavour Standard Model, NSI,
the 3+1 sterile extension, or any combination, with $H_j$ simply changing
size ($3\times3$ or $4\times4$) and content accordingly; nothing in this
module branches on which scenario is active. In particular, an optional
neutron density $n_n(\bar x_j)$ sampled alongside $n_e(\bar x_j)$ is
forwarded straight through to \pyfunc{hamiltonian\_flavour} as
\code{n\_n\_mol\_cm3}, enabling the 3+1 sterile extension's neutral-current
term (\cref{ch:core-common}) with no change to this module's own logic:
omitting it (the default) reproduces the pre-existing CC-only behaviour
exactly.

\subsection*{Composing the total evolutor}
The $N$ per-segment evolutors are composed into the total operator via the
ordered product already established generically in
\cref{ch:core-common}:
\begin{equation}
  S(x_N,x_0) = S_{N-1}\,S_{N-2}\cdots S_1\,S_0,
\end{equation}
using \pyfunc{compose\_segment\_evolutors} (\cref{ch:core-common}) with its
binary-tree reduction for $O(\log N)$ parallel depth on GPU rather than
$O(N)$ sequential steps. Optionally, \pyfunc{evolutor\_numerical} can
instead return the \emph{history} of partial products
$S(x_j,x_0)$ for every $j=0,\dots,N$ (an accumulating loop rather than the
tree reduction, since every intermediate result is needed): useful for
path-resolved quantities, e.g.\ probabilities evaluated at every crossed
layer of a medium rather than only at the final boundary.
\end{theorybox}

\begin{implbox}
\begin{itemize}
  \item \pyfunc{evolutor\_numerical\_segment(oscillation, E\_MeV, n\_e\_mol\_cm3,
    dx\_evolution, *, n\_n\_mol\_cm3=None, ...)}: builds one $H_j$ per segment
    (the segment dimension is the final axis of
    \code{n\_e\_mol\_cm3}/\code{dx\_evolution}), broadcasting
    \code{oscillation.mass\_spectrum} and \code{antinu} to the segment shape
    via \code{dataclasses.replace} before calling \pyfunc{hamiltonian\_flavour}
    once for the whole segment batch, then exponentiates:
    \code{torch.linalg.matrix\_exp(-1j * H * dx[...,None,None])}.
  \item \pyfunc{evolutor\_numerical(oscillation, E\_MeV, n\_e\_mol\_cm3,
    dx\_evolution, *, n\_n\_mol\_cm3=None, return\_history=False, ...)}: calls
    the above, then either composes the segments into a single operator
    (\pyfunc{compose\_segment\_evolutors}) or accumulates and returns the
    full per-segment history.
  \item Both functions run under \code{@torch.no\_grad()}: the numerical
    evolutor is used for forward propagation, not for parameter inference
    through automatic differentiation.
\end{itemize}
\end{implbox}

\begin{notebox}
\textbf{Accuracy is entirely a discretisation choice, made outside this
module.} \texttt{core.numerical} performs no adaptive step-size control and
no local error estimate: the accuracy of $S(x_N,x_0)$ is fixed entirely by
how finely the \pyclass{Trajectory} passed in discretises the path (how
many segments $N$, and which quadrature rule fixed $\bar x_j$, previous
section) -- both decided by the calling medium-specific module before
\texttt{core.numerical} ever sees the problem. Refining the trajectory
(more, smaller segments) is the only way to improve accuracy; there is no
mechanism to detect that a given $N$ is insufficient. This is the
computational cost paid for generality: unlike the perturbative
approximation (\cref{ch:core-perturbative}), which is only valid in a
restricted regime but is $\mathcal{O}(1)$ in the number of density
evaluations, the numerical method is valid everywhere but its accuracy is
$\mathcal{O}(N)$ in the number of \pyfunc{hamiltonian\_flavour}
evaluations and matrix exponentials.
\end{notebox}
